\documentclass[10pt,letterpaper]{article}
\usepackage[utf8]{inputenc}
\usepackage[T1]{fontenc}
\usepackage{mathptmx}
\renewcommand{\ttdefault}{txtt}
\usepackage[left=1.9cm,right=1.9cm,top=3.2cm,bottom=2.5cm,headheight=1.7cm,headsep=0.4cm,footskip=0.9cm]{geometry}
\usepackage{xcolor}
\usepackage{graphicx}
\usepackage[font={footnotesize,it},labelsep=colon,hypcap=false]{caption}
\usepackage{fancyhdr}
\usepackage{eso-pic}
\usepackage{titlesec}
\usepackage{listings}
\usepackage[hidelinks]{hyperref}

\linespread{1.1}
\setlength{\parindent}{0pt}
\setlength{\parskip}{6pt}
\definecolor{gris}{gray}{0.35}
\setlength{\fboxsep}{0pt}
\setlength{\fboxrule}{0.4pt}
\captionsetup[figure]{name=Imagen}
\titleformat*{\section}{\normalsize\bfseries}
\titlespacing*{\section}{0pt}{14pt}{6pt}

\pagestyle{fancy}
\fancyhf{}
\renewcommand{\headrulewidth}{0pt}
\fancyhead[L]{\includegraphics[height=1.4cm]{logo_usm}}
\fancyhead[R]{\footnotesize UNIVERSIDAD TÉCNICA FEDERICO SANTA MARÍA\\DEPARTAMENTO DE INDUSTRIAS\\CAMPUS SANTIAGO VITACURA\\SEGUNDO SEMESTRE 2026}
\fancyfoot[R]{\footnotesize\thepage}

\AddToShipoutPictureBG{\AtPageLowerLeft{\setlength{\unitlength}{1bp}\color[gray]{0.6}\put(24,24){\framebox(564,744){}}}}

\definecolor{clave}{HTML}{1A8BCB}
\definecolor{funcion}{HTML}{A8A898}
\definecolor{texto}{HTML}{E08A14}
\definecolor{alias}{HTML}{A0422E}
\definecolor{codigo}{HTML}{212121}

\makeatletter
\newcommand{\numero}[1]{\ifnum\lst@mode=\lst@Pmode\textcolor{texto}{#1}\else#1\fi}
\makeatother

\lstset{
  morestring=[b]',
  moredelim=[s][\color{alias}]{`}{`},
  keywords={SELECT,FROM,WHERE,AND,OR,NOT,IS,NULL,IN,BETWEEN,ORDER,BY,GROUP,HAVING,LIMIT,DESC,ASC,AS,JOIN,ON,DISTINCT,SET,UPDATE,DELETE,INSERT,INTO,VALUES,AVG},
  keywords=[2]{COUNT,SUM,MIN,MAX},
  basicstyle=\ttfamily\footnotesize\color{codigo},
  keywordstyle=\bfseries\color{clave},
  keywordstyle=[2]\color{funcion},
  stringstyle=\color{texto},
  frame=tb,
  columns=fullflexible,
  keepspaces=true,
  upquote=true,
  showstringspaces=false,
  literate={á}{{\'a}}1 {é}{{\'e}}1 {í}{{\'i}}1 {ó}{{\'o}}1 {ú}{{\'u}}1 {ñ}{{\~n}}1
           {0}{{\numero{0}}}1 {1}{{\numero{1}}}1 {2}{{\numero{2}}}1 {3}{{\numero{3}}}1 {4}{{\numero{4}}}1 {5}{{\numero{5}}}1 {6}{{\numero{6}}}1 {7}{{\numero{7}}}1 {8}{{\numero{8}}}1 {9}{{\numero{9}}}1
}

\begin{document}

\begin{center}
\vspace*{4cm}
{\fontsize{26}{32}\selectfont\bfseries\textsc{Laboratorio 4}\par}
\vspace{10pt}
{\color{lightgray}\rule{0.75\linewidth}{0.5pt}}\par
\vspace{2pt}
{\itshape Entregable SQL (1/2): base de datos Austral Editores\par}
{\color{lightgray}\rule{0.75\linewidth}{0.5pt}}\par
\vspace{12pt}
{\fontsize{22}{28}\selectfont\bfseries\textsc{Sistemas de información para la gestión.}\par}
\vspace{4.5cm}
{\color{gris}\itshape Daniela Román Schuffenegger, rol 202160649-1, \href{mailto:daniela.roman@usm.cl}{\underline{daniela.roman@usm.cl}}\\
Paralelo 101\par}
\vspace{1.8cm}
{\color{gris}\itshape Profesor: José Luis Sáez.\\
Ayudantes: Valeria Bobadilla, Catalina Bodaleo, Luis Contreras, Benjamín Jara.\\
Fecha: 31 de octubre de 2026\par}
\end{center}
\vspace{0.8cm}
\textbf{Repositorio Github: \href{https://github.com/DanielaRoman7/ICN292-Lab4-Roman-Daniela}{https://github.com/DanielaRoman7/ICN292-Lab4-Roman-Daniela}}
\newpage

\section*{Consulta 1: autores que no son de Chile}

\begin{minipage}{\linewidth}
\begin{lstlisting}
SELECT nombre AS `Nombre`, genero AS `Género`
FROM autores
WHERE pais <> 'Chile' OR pais IS NULL
ORDER BY nombre;
\end{lstlisting}
\centering
\fcolorbox{gray}{white}{\includegraphics[width=0.25\linewidth]{capturas/consulta01}}
\captionof{figure}{Autores que no son de Chile}
\end{minipage}

\bigskip

Para Kenji Arakawa, \lstinline|pais <> 'Chile'| no da verdadero sino \lstinline|NULL|, porque su país está vacío y cualquier comparación con \lstinline|NULL| da \lstinline|NULL|; como el \lstinline|WHERE| solo deja pasar las filas verdaderas, Kenji quedaba fuera, y por eso agregué \lstinline|OR pais IS NULL|.

\section*{Consulta 2: libros entre 190 y 220 páginas}

\begin{minipage}{\linewidth}
\begin{lstlisting}
SELECT titulo AS `Título`, paginas AS `Páginas`
FROM libros
WHERE paginas BETWEEN 190 AND 220
ORDER BY paginas;
\end{lstlisting}
\centering
\fcolorbox{gray}{white}{\includegraphics[width=0.25\linewidth]{capturas/consulta02}}
\captionof{figure}{Libros entre 190 y 220 páginas}
\end{minipage}

\section*{Consulta 3: resumen del catálogo}

\begin{minipage}{\linewidth}
\begin{lstlisting}
SELECT COUNT(*) AS `Total de libros`,
       SUM(ejemplares_vendidos) AS `Total de ejemplares`,
       AVG(paginas) AS `Promedio de páginas`,
       MIN(anio) AS `Año más antiguo`,
       MAX(ejemplares_vendidos) AS `Máximo por libro`
FROM libros;
\end{lstlisting}
\centering
\fcolorbox{gray}{white}{\includegraphics[width=0.7\linewidth]{capturas/consulta03}}
\captionof{figure}{Resumen del catálogo}
\end{minipage}

\section*{Consulta 4: géneros de los autores}

\begin{minipage}{\linewidth}
\begin{lstlisting}
SELECT DISTINCT genero AS `Género`
FROM autores
ORDER BY genero;
\end{lstlisting}
\centering
\fcolorbox{gray}{white}{\includegraphics[width=0.1\linewidth]{capturas/consulta04}}
\captionof{figure}{Géneros de los autores}
\end{minipage}

\section*{Consulta 5: autores con más de un libro en 2024}

\begin{minipage}{\linewidth}
\begin{lstlisting}
SELECT a.nombre AS `Nombre`, COUNT(*) AS `Libros en 2024`
FROM autores a
JOIN libros l ON a.id_autor = l.id_autor
WHERE l.anio = 2024
GROUP BY a.id_autor, a.nombre
HAVING COUNT(*) > 1
ORDER BY a.nombre;
\end{lstlisting}
\centering
\fcolorbox{gray}{white}{\includegraphics[width=0.3\linewidth]{capturas/consulta05}}
\captionof{figure}{Autores con más de un libro en 2024}
\end{minipage}

\section*{Consulta 6: los libros de 2024 más vendidos}

\begin{minipage}{\linewidth}
\begin{lstlisting}
SELECT titulo AS `Título`, ejemplares_vendidos AS `Ejemplares vendidos`
FROM libros
WHERE anio = 2024 AND ejemplares_vendidos > 2000
ORDER BY ejemplares_vendidos DESC
LIMIT 3;
\end{lstlisting}
\centering
\fcolorbox{gray}{white}{\includegraphics[width=0.35\linewidth]{capturas/consulta06}}
\captionof{figure}{Libros de 2024 con más de 2.000 ejemplares}
\end{minipage}

\section*{Consulta 7: agregar y borrar «Papel Mojado»}

\begin{minipage}{\linewidth}
\textbf{Agrego el libro}
\begin{lstlisting}
INSERT INTO libros (id_autor, titulo, paginas, ejemplares_vendidos, anio)
VALUES (3, 'Papel Mojado', 260, 50, 2024);
\end{lstlisting}
\centering
\fcolorbox{gray}{white}{\includegraphics[width=0.85\linewidth]{capturas/consulta07a}}
\captionof{figure}{INSERT de Papel Mojado}
\end{minipage}

\bigskip

\begin{minipage}{\linewidth}
\textbf{Reviso que quedó guardado}
\begin{lstlisting}
SELECT id_libro AS `ID libro`,
       id_autor AS `ID autor`,
       id_coleccion AS `ID colección`,
       titulo AS `Título`,
       paginas AS `Páginas`,
       ejemplares_vendidos AS `Ejemplares vendidos`,
       anio AS `Año`
FROM libros
WHERE titulo = 'Papel Mojado';
\end{lstlisting}
\centering
\fcolorbox{gray}{white}{\includegraphics[width=0.75\linewidth]{capturas/consulta07b}}
\captionof{figure}{Papel Mojado guardado en la tabla}
\end{minipage}

\bigskip

\begin{minipage}{\linewidth}
\textbf{Lo borro filtrando por el título}
\begin{lstlisting}
DELETE FROM libros
WHERE titulo = 'Papel Mojado';
\end{lstlisting}
\centering
\fcolorbox{gray}{white}{\includegraphics[width=0.85\linewidth]{capturas/consulta07c}}
\captionof{figure}{DELETE de Papel Mojado}
\end{minipage}

\bigskip

\begin{minipage}{\linewidth}
\textbf{Vuelvo a buscarlo para confirmar que ya no está}
\begin{lstlisting}
SELECT id_libro AS `ID libro`,
       id_autor AS `ID autor`,
       id_coleccion AS `ID colección`,
       titulo AS `Título`,
       paginas AS `Páginas`,
       ejemplares_vendidos AS `Ejemplares vendidos`,
       anio AS `Año`
FROM libros
WHERE titulo = 'Papel Mojado';
\end{lstlisting}
\centering
\fcolorbox{gray}{white}{\includegraphics[width=0.7\linewidth]{capturas/consulta07d}}
\captionof{figure}{Papel Mojado ya no está en la tabla}
\end{minipage}

\bigskip

\begin{minipage}{\linewidth}
\textbf{Cuento cuántos libros quedan con ese título}
\begin{lstlisting}
SELECT COUNT(*) AS `Libros encontrados`
FROM libros
WHERE titulo = 'Papel Mojado';
\end{lstlisting}
\centering
\fcolorbox{gray}{white}{\includegraphics[width=0.15\linewidth]{capturas/consulta07e}}
\captionof{figure}{Conteo de libros llamados Papel Mojado}
\end{minipage}

\section*{Consulta 8: sumar 100 ejemplares a «Cielo de Enero»}

\begin{minipage}{\linewidth}
\textbf{Cómo estaba antes del cambio}
\begin{lstlisting}
SELECT titulo AS `Título`, ejemplares_vendidos AS `Ejemplares vendidos`
FROM libros
WHERE titulo = 'Cielo de Enero';
\end{lstlisting}
\centering
\fcolorbox{gray}{white}{\includegraphics[width=0.3\linewidth]{capturas/consulta08a}}
\captionof{figure}{Cielo de Enero antes del UPDATE}
\end{minipage}

\bigskip

\begin{minipage}{\linewidth}
\textbf{El UPDATE, que suma sobre el valor actual}
\begin{lstlisting}
UPDATE libros
SET ejemplares_vendidos = ejemplares_vendidos + 100
WHERE titulo = 'Cielo de Enero';
\end{lstlisting}
\centering
\fcolorbox{gray}{white}{\includegraphics[width=0.85\linewidth]{capturas/consulta08b}}
\captionof{figure}{UPDATE de Cielo de Enero}
\end{minipage}

\bigskip

\begin{minipage}{\linewidth}
\textbf{Cómo quedó después}
\begin{lstlisting}
SELECT titulo AS `Título`, ejemplares_vendidos AS `Ejemplares vendidos`
FROM libros
WHERE titulo = 'Cielo de Enero';
\end{lstlisting}
\centering
\fcolorbox{gray}{white}{\includegraphics[width=0.3\linewidth]{capturas/consulta08c}}
\captionof{figure}{Cielo de Enero después del UPDATE}
\end{minipage}

\section*{Consulta 9: libros y sus colecciones}

\begin{minipage}{\linewidth}
\textbf{Cargo tres colecciones de prueba}
\begin{lstlisting}
INSERT INTO colecciones (titulo, fecha_publicacion) VALUES
('Narrativa Austral', '2024-11-15'),
('Ideas y Versos', '2024-08-20'),
('Pequeños Lectores', '2025-03-10');
\end{lstlisting}
\centering
\fcolorbox{gray}{white}{\includegraphics[width=0.85\linewidth]{capturas/consulta09a}}
\captionof{figure}{INSERT de las colecciones}
\end{minipage}

\bigskip

\begin{minipage}{\linewidth}
\textbf{Asigno las novelas de Elena Pizarro y Nora Keller a la colección 1}
\begin{lstlisting}
UPDATE libros
SET id_coleccion = 1
WHERE id_autor IN (1, 5);
\end{lstlisting}
\centering
\fcolorbox{gray}{white}{\includegraphics[width=0.85\linewidth]{capturas/consulta09b}}
\captionof{figure}{Libros asignados a la colección 1}
\end{minipage}

\bigskip

\begin{minipage}{\linewidth}
\textbf{Asigno los libros de Harold Finch y Rodrigo Bahamonde a la colección 2}
\begin{lstlisting}
UPDATE libros
SET id_coleccion = 2
WHERE id_autor IN (2, 3);
\end{lstlisting}
\centering
\fcolorbox{gray}{white}{\includegraphics[width=0.85\linewidth]{capturas/consulta09c}}
\captionof{figure}{Libros asignados a la colección 2}
\end{minipage}

\bigskip

\begin{minipage}{\linewidth}
\textbf{Asigno los libros de Valentina Correa a la colección 3}
\begin{lstlisting}
UPDATE libros
SET id_coleccion = 3
WHERE id_autor = 4;
\end{lstlisting}
\centering
\fcolorbox{gray}{white}{\includegraphics[width=0.85\linewidth]{capturas/consulta09d}}
\captionof{figure}{Libros asignados a la colección 3}
\end{minipage}

\bigskip

\begin{minipage}{\linewidth}
\textbf{Listo cada libro con su colección y la fecha de publicación}
\begin{lstlisting}
SELECT l.titulo AS `Libro`,
       c.titulo AS `Colección`,
       c.fecha_publicacion AS `Fecha de publicación`
FROM libros l
JOIN colecciones c ON l.id_coleccion = c.id_coleccion
ORDER BY c.fecha_publicacion, l.titulo;
\end{lstlisting}
\centering
\fcolorbox{gray}{white}{\includegraphics[width=0.5\linewidth]{capturas/consulta09e}}
\captionof{figure}{Libros con su colección}
\end{minipage}

\section*{Consulta 10: pagos a los autores}

\begin{minipage}{\linewidth}
\textbf{Cargo cuatro pagos de prueba}
\begin{lstlisting}
INSERT INTO pagos (id_autor, fecha_pago, monto) VALUES
(1, '2026-03-31', 450000.00),
(4, '2026-03-31', 1250000.00),
(1, '2026-06-30', 385500.00),
(5, '2026-06-30', 980000.00);
\end{lstlisting}
\centering
\fcolorbox{gray}{white}{\includegraphics[width=0.85\linewidth]{capturas/consulta10a}}
\captionof{figure}{INSERT de los pagos}
\end{minipage}

\bigskip

\begin{minipage}{\linewidth}
\textbf{Primera sentencia: cada pago con su autor}
\begin{lstlisting}
SELECT a.nombre AS `Nombre`,
       p.fecha_pago AS `Fecha de pago`,
       p.monto AS `Monto`
FROM pagos p
JOIN autores a ON p.id_autor = a.id_autor
ORDER BY p.fecha_pago, a.nombre;
\end{lstlisting}
\centering
\fcolorbox{gray}{white}{\includegraphics[width=0.4\linewidth]{capturas/consulta10b}}
\captionof{figure}{Pagos con el nombre del autor}
\end{minipage}

\bigskip

\begin{minipage}{\linewidth}
\textbf{Segunda sentencia: total pagado a cada autor}
\begin{lstlisting}
SELECT a.nombre AS `Nombre`, SUM(p.monto) AS `Total pagado`
FROM autores a
JOIN pagos p ON a.id_autor = p.id_autor
GROUP BY a.id_autor, a.nombre
ORDER BY a.nombre;
\end{lstlisting}
\centering
\fcolorbox{gray}{white}{\includegraphics[width=0.25\linewidth]{capturas/consulta10c}}
\captionof{figure}{Total pagado a cada autor}
\end{minipage}

\end{document}
